\section{Metodología}

\subsection{Gestión del código fuente}
El código fuente del proyecto se gestionará usando un sistema de control de versiones Git, permitiendo mantener el historial de cambios y asegurar el orden y la claridad en el desarrollo.

\subsection{Política de ramas}
Se seguirá un flujo de trabajo sencillo:

\begin{itemize}
    \item \textbf{main}: versión estable del proyecto.
    \item \textbf{develop}: rama de integración de nuevas funcionalidades antes de su fusión en \textit{main}.
    \item \textbf{feature/}: ramas derivadas de \textit{develop}, donde se desarrollará cada nueva funcionalidad de forma aislada.
    \item \textbf{hotfix/}: ramas destinadas a la corrección de errores críticos directamente sobre \textit{main}.
\end{itemize}

\subsection{Política de versiones del software a desarrollar}
El software seguirá un esquema de versionado semántico de tres niveles (X.Y.Z):

\begin{itemize}
    \item \textbf{X}: versión mayor, aplicada ante cambios incompatibles o rediseños importantes del sistema.
    \item \textbf{Y}: versión menor, correspondiente a la incorporación de nuevas funcionalidades.
    \item \textbf{Z}: versión de parche, destinada a la corrección de errores o problemas menores.
\end{itemize}

\subsection{Política de mensajes de commit}
Se seguirá el estándar de \textit{Conventional Commits} para mantener un historial claro y estructurado:

\begin{center}
\texttt{<tipo>[alcance opcional]: <descripción corta>}
\end{center}

Tipos de commits:

\begin{itemize}
    \item \textbf{feat}: nueva funcionalidad.
    \item \textbf{fix}: corrección de errores.
    \item \textbf{docs}: cambios en la documentación.
    \item \textbf{style}: cambios de formato que no afectan a la funcionalidad.
    \item \textbf{refactor}: cambios en el código que no alteran el comportamiento.
    \item \textbf{test}: adición o modificación de pruebas.
    \item \textbf{chore}: tareas de mantenimiento o configuración.
\end{itemize}

\subsection{Estimación de historias de usuario: tallas de camiseta}
Para estimar el esfuerzo de las historias de usuario se utiliza la técnica de tallas de camiseta (\textit{T-shirt sizing}), que asigna un tamaño relativo a cada historia: XS (extra pequeña), S (pequeña), M (mediana), L (grande) y XL (extra grande).

La talla se define en función de la complejidad, el riesgo y la cantidad de trabajo, comparando cada historia con otras ya conocidas. En sesiones de refinamiento, el equipo acuerda la talla por consenso, sin necesidad de estimaciones en horas exactas.

Esta técnica permite priorizar el backlog combinando valor (prioridad) y esfuerzo (talla), facilitando la planificación ágil y la selección de historias por sprint.

\subsection{Gestión del proyecto}
El proyecto se gestionará con una adaptación ligera de Scrum, ajustada al contexto de un Trabajo de Fin de Grado (TFG), donde el equipo es reducido y el ritmo de trabajo es más flexible que en un entorno profesional.

\begin{itemize}
    \item \textbf{No se realizan daily standups}: debido al tamaño del equipo y la naturaleza del TFG, no se llevan a cabo reuniones diarias de 15 minutos. En su lugar, la comunicación se mantiene de forma continua.
    
    \item \textbf{Review y planning fusionados}: las reuniones de revisión (\textit{review}) y planificación (\textit{sprint planning}) se combinan en una única sesión por sprint. En esta sesión se revisa lo completado, se evalúa el avance respecto a los objetivos del sprint y se planifica el siguiente.

    \item \textbf{Sprints con objetivos definidos}: se mantiene la estructura de sprints, cada uno con un objetivo claro y un conjunto de historias de usuario o tareas a completar.
\end{itemize}